\documentclass[10pt,a4paper]{amsart}
\usepackage{amsmath,amssymb,amsthm,graphicx}
\usepackage{verbatim}
\newenvironment{reference}{\comment}{\endcomment}
%\newenvironment{reference}{}{}
\newenvironment{slogan}{\comment}{\endcomment}
\newenvironment{history}{\comment}{\endcomment}

% For commutative diagrams we use Xy-pic
\usepackage[all]{xy}

% We use 2cell for 2-commutative diagrams.
\xyoption{2cell}
% Theorem environments.
%
\theoremstyle{plain}
\newtheorem{theorem}[subsection]{Theorem}
\newtheorem{proposition}[subsection]{Proposition}
\newtheorem{lemma}[subsection]{Lemma}

\theoremstyle{definition}
\newtheorem{definition}[subsection]{Definition}
\newtheorem{example}[subsection]{Example}
\newtheorem{exercise}[subsection]{Exercise}
\newtheorem{situation}[subsection]{Situation}

\theoremstyle{remark}
\newtheorem{remark}[subsection]{Remark}
\newtheorem{remarks}[subsection]{Remarks}

\numberwithin{equation}{subsection}

% Macros
%
\def\lim{\mathop{\mathrm{lim}}\nolimits}
\def\colim{\mathop{\mathrm{colim}}\nolimits}
\def\Spec{\mathop{\mathrm{Spec}}}
\def\Hom{\mathop{\mathrm{Hom}}\nolimits}
\def\Ext{\mathop{\mathrm{Ext}}\nolimits}
\def\SheafHom{\mathop{\mathcal{H}\!\mathit{om}}\nolimits}
\def\SheafExt{\mathop{\mathcal{E}\!\mathit{xt}}\nolimits}
\def\Sch{\mathit{Sch}}
\def\Mor{\mathop{\mathrm{Mor}}\nolimits}
\def\Ob{\mathop{\mathrm{Ob}}\nolimits}
\def\Sh{\mathop{\mathit{Sh}}\nolimits}
\def\NL{\mathop{N\!L}\nolimits}
\def\CH{\mathop{\mathrm{CH}}\nolimits}
\def\proetale{{pro\text{-}\acute{e}tale}}
\def\etale{{\acute{e}tale}}
\def\QCoh{\mathit{QCoh}}
\def\Ker{\mathop{\mathrm{Ker}}}
\def\Im{\mathop{\mathrm{Im}}}
\def\Coker{\mathop{\mathrm{Coker}}}
\def\Coim{\mathop{\mathrm{Coim}}}

% Boxtimes
%
\DeclareMathSymbol{\boxtimes}{\mathbin}{AMSa}{"02}

%
% Macros for moduli stacks/spaces
%
\def\QCohstack{\mathcal{QC}\!\mathit{oh}}
\def\Cohstack{\mathcal{C}\!\mathit{oh}}
\def\Spacesstack{\mathcal{S}\!\mathit{paces}}
\def\Quotfunctor{\mathrm{Quot}}
\def\Hilbfunctor{\mathrm{Hilb}}
\def\Curvesstack{\mathcal{C}\!\mathit{urves}}
\def\Polarizedstack{\mathcal{P}\!\mathit{olarized}}
\def\Complexesstack{\mathcal{C}\!\mathit{omplexes}}
% \Pic is the operator that assigns to X its picard group, usage \Pic(X)
% \Picardstack_{X/B} denotes the Picard stack of X over B
% \Picardfunctor_{X/B} denotes the Picard functor of X over B
\def\Pic{\mathop{\mathrm{Pic}}\nolimits}
\def\Picardstack{\mathcal{P}\!\mathit{ic}}
\def\Picardfunctor{\mathrm{Pic}}
\def\Deformationcategory{\mathcal{D}\!\mathit{ef}}
\usepackage{cleveref}
\crefname{originalreference}{original reference}{original references}
\Crefname{originalreference}{Original reference}{Original references}
\makeatletter
\newlabel{remark-integral}{{0EDV}{source}{}{sourcepointer-1}{}}
\newlabel{remark-integral@cref}{{[originalreference][0][]0EDV}{[1][0][]source}{}{}{}}
\makeatother
\crefname{originalreference}{original reference}{original references}
\Crefname{originalreference}{Original reference}{Original references}
\makeatletter
\newlabel{situation-setup}{{02QL}{source}{}{sourcepointer-2}{}}
\newlabel{situation-setup@cref}{{[originalreference][0][]02QL}{[1][0][]source}{}{}{}}
\makeatother
\crefname{originallemma}{original lemma}{original lemmas}
\Crefname{originallemma}{Original lemma}{Original lemmas}
\makeatletter
\newlabel{spaces-morphisms-lemma-base-change-proper}{{04WP}{source}{}{sourcepointer-3}{}}
\newlabel{spaces-morphisms-lemma-base-change-proper@cref}{{[originallemma][0][]04WP}{[1][0][]source}{}{}{}}
\makeatother
\crefname{originallemma}{original lemma}{original lemmas}
\Crefname{originallemma}{Original lemma}{Original lemmas}
\makeatletter
\newlabel{spaces-morphisms-lemma-closed-immersion-proper}{{04XZ}{source}{}{sourcepointer-4}{}}
\newlabel{spaces-morphisms-lemma-closed-immersion-proper@cref}{{[originallemma][0][]04XZ}{[1][0][]source}{}{}{}}
\makeatother
\crefname{originallemma}{original lemma}{original lemmas}
\Crefname{originallemma}{Original lemma}{Original lemmas}
\makeatletter
\newlabel{spaces-morphisms-lemma-universally-closed-permanence}{{04NX}{source}{}{sourcepointer-5}{}}
\newlabel{spaces-morphisms-lemma-universally-closed-permanence@cref}{{[originallemma][0][]04NX}{[1][0][]source}{}{}{}}
\makeatother
\crefname{originallemma}{original lemma}{original lemmas}
\Crefname{originallemma}{Original lemma}{Original lemmas}
\makeatletter
\newlabel{spaces-properties-lemma-map-into-reduction}{{03JJ}{source}{}{sourcepointer-6}{}}
\newlabel{spaces-properties-lemma-map-into-reduction@cref}{{[originallemma][0][]03JJ}{[1][0][]source}{}{}{}}
\makeatother
\crefname{originallemma}{original lemma}{original lemmas}
\Crefname{originallemma}{Original lemma}{Original lemmas}
\makeatletter
\newlabel{spaces-properties-lemma-points-cartesian}{{03H4}{source}{}{sourcepointer-7}{}}
\newlabel{spaces-properties-lemma-points-cartesian@cref}{{[originallemma][0][]03H4}{[1][0][]source}{}{}{}}
\makeatother
\title{Every integral closed subspace has a unique dominant image, proper and surjective when the morphism is proper}
\author{The Stacks Project Authors}
\date{stacks-project:spaces-chow.tex:local-snapshot-20261012; GNU FDL 1.2 or later}
\begin{document}
\maketitle
\noindent\textit{Editorial note.} All delivered mathematical text is the author's own, copied from the pinned Stacks Project chapter file spaces-chow.tex; no statement or proof step was rewritten, reordered or removed.\par\medskip
\noindent\textit{Editorial note.} The counted proof is the author's complete argument: a construction of the image through the generic point with a reduction argument for the morphism, and a short properness and surjectivity chain in the proper case\par\medskip
\noindent\textit{Editorial note.} The label lemma-proper-image names the counted statement, and the proof environment immediately following it in the source is delivered verbatim as its proof.\par\medskip
\noindent\textit{Editorial note.} The counted statement is carried by the author's own lemma environment in the chapter source.\par\medskip
\noindent\textit{Allowed standard prerequisites (editorial).}\begin{itemize}
\item Integral closed subspaces and their generic points, the closure of the image of a generic point being the generic point of an integral closed subspace, and the reduction of a closed subspace
\item The cartesian description of the points of the preimage of a closed subspace, and the universal property of the reduction
\item Stability of properness under composition, under base change and for closed immersions, and closedness of the image of a proper morphism
\end{itemize}\medskip
\ifdefined\hypertarget\hypertarget{sourcepointer-1}{}\fi
\noindent\textit{Source reference.} Original reference 0EDV: sources/stacks-spaces-chow/pinned-original-source.tex, source lines 311--311. Prerequisite pointer invoked at proof line(s) 838; it names a standard fact of the same chapter that the counted argument uses, and every step of the counted proof built on it is retained verbatim in the delivered body.\par\medskip
\ifdefined\hypertarget\hypertarget{sourcepointer-2}{}\fi
\noindent\textit{Source reference.} Original reference 02QL: sources/stacks-spaces-chow/pinned-original-source.tex, source lines 264--264. Prerequisite pointer invoked at proof line(s) 821; it names a standard fact of the same chapter that the counted argument uses, and every step of the counted proof built on it is retained verbatim in the delivered body.\par\medskip
\ifdefined\hypertarget\hypertarget{sourcepointer-3}{}\fi
\noindent\textit{Source reference.} Original lemma 04WP: sources/stacks-chow/tags.txt, source lines 6367--6367. Prerequisite pointer invoked at proof line(s) 848; it names a standard fact of a sibling chapter that the counted argument uses, and every step of the counted proof built on it is retained verbatim in the delivered body.\par\medskip
\ifdefined\hypertarget\hypertarget{sourcepointer-4}{}\fi
\noindent\textit{Source reference.} Original lemma 04XZ: sources/stacks-chow/tags.txt, source lines 6412--6412. Prerequisite pointer invoked at proof line(s) 849; it names a standard fact of a sibling chapter that the counted argument uses, and every step of the counted proof built on it is retained verbatim in the delivered body.\par\medskip
\ifdefined\hypertarget\hypertarget{sourcepointer-5}{}\fi
\noindent\textit{Source reference.} Original lemma 04NX: sources/stacks-chow/tags.txt, source lines 6095--6095. Prerequisite pointer invoked at proof line(s) 852; it names a standard fact of a sibling chapter that the counted argument uses, and every step of the counted proof built on it is retained verbatim in the delivered body.\par\medskip
\ifdefined\hypertarget\hypertarget{sourcepointer-6}{}\fi
\noindent\textit{Source reference.} Original lemma 03JJ: sources/stacks-chow/tags.txt, source lines 4641--4641. Prerequisite pointer invoked at proof line(s) 842; it names a standard fact of a sibling chapter that the counted argument uses, and every step of the counted proof built on it is retained verbatim in the delivered body.\par\medskip
\ifdefined\hypertarget\hypertarget{sourcepointer-7}{}\fi
\noindent\textit{Source reference.} Original lemma 03H4: sources/stacks-chow/tags.txt, source lines 4553--4553. Prerequisite pointer invoked at proof line(s) 839; it names a standard fact of a sibling chapter that the counted argument uses, and every step of the counted proof built on it is retained verbatim in the delivered body.\par\medskip
\begin{lemma}
\label{lemma-proper-image}
In Situation \ref{situation-setup} let $X,Y/B$ be good and let $f : X \to Y$
be a morphism over $B$. If $Z \subset X$ is an integral closed subspace, then
there exists a unique integral closed subspace $Z' \subset Y$ such that there
is a commutative diagram
$$
\xymatrix{
Z \ar[r] \ar[d] & X \ar[d]^f \\
Z' \ar[r] & Y
}
$$
with $Z \to Z'$ dominant. If $f$ is proper, then $Z \to Z'$ is proper
and surjective.
\end{lemma}

\begin{proof}
Let $\xi \in |Z|$ be the generic point. Let $Z' \subset Y$ be the integral
closed subspace whose generic point is $\xi' = f(\xi)$, see
Remark \ref{remark-integral}. Since $\xi \in |f^{-1}(Z')| = |f|^{-1}(|Z'|)$
by Properties of Spaces, Lemma \ref{spaces-properties-lemma-points-cartesian}
and since $Z$ is the reduced with $|Z| = \overline{\{\xi\}}$
we see that $Z \subset f^{-1}(Z')$ as closed subspaces of $X$ (see
Properties of Spaces, Lemma \ref{spaces-properties-lemma-map-into-reduction}).
Thus we obtain our morphism $Z \to Z'$.
This morphism is dominant as the generic point of $Z$
maps to the generic point of $Z'$. Uniqueness of $Z'$ is clear.
If $f$ is proper, then $Z \to Y$ is proper as a composition
of proper morphisms (Morphisms of Spaces, Lemmas
\ref{spaces-morphisms-lemma-base-change-proper} and
\ref{spaces-morphisms-lemma-closed-immersion-proper}).
Then we conclude that $Z \to Z'$ is proper by
Morphisms of Spaces, Lemma
\ref{spaces-morphisms-lemma-universally-closed-permanence}.
Surjectivity then follows as the image of a proper morphism is closed.
\end{proof}

\end{document}
